% Mechanical relocation generated from the preservation ledger.
% Existing prose is included byte-for-byte from preserved blocks.

\section{Current Limitations}
\subsection{Large OD Layouts}

An explicit dense measurement-by-OD matrix is infeasible for a full network,
and repeated detailed graph traversals can dominate runtime even when only a
small demand block changes. Matrix-free and sharded representations bound
storage but do not remove the cost of an unnecessarily high-dimensional demand
model. This motivates the reduced gravity parameterization and the separation
between response construction, estimation, and detailed post-fit validation.

\subsection{Identification and Model Dependence}

Stop counts generally identify combinations of OD flows rather than every
cell. Priors, gravity structure, production constraints, entropy, or seeds can
select stable estimates, but they do not create observational information.
Reported uncertainty and sensitivity must distinguish data-driven precision
from restrictions imposed by the demand model.

\subsection{Fixed Routing and Capacity}

Routing probabilities are conditional on a fixed timetable, generalized costs,
destination rules, and dispersion. Demand does not change travel time,
capacity, denied boarding, or route utility. This is internally consistent with
the current assignment but excludes crowding feedback and capacity-constrained
equilibria. In heavily loaded conditions, a capacity-aware outer assignment and
re-estimation loop would be required.

\subsection{Observation and Data Limitations}

The count likelihood assumes a correct mapping from operational records to
timetable events and a suitable conditional count distribution. Missing
records are not zeros, transfer events are not journey origins or final
destinations, and systematic detector failures cannot be repaired by the OD
model. Temporal aggregation and physical-stop mappings can remove or create
apparent distinctions, so their uncertainty remains a source of model risk.

\subsection{Reduced Journey Choices and External Inputs}

The reduced model is conditional on feasible alternatives, fixed shares,
travel-time and transfer summaries, attractiveness values, and journey
productions or a production basis. Transfer limits and choice-set pruning may
exclude behavior that occurs in practice. Gravity parameters can diagnose weak
information within this specification, but optimal placement of new
measurements is not yet implemented.

\subsection{Inference and Computational Limits}

MAP provides a mode and local curvature, not a complete uncertainty
distribution. Variational inference provides an approximate posterior whose
quality depends on the guide and optimization. Full cell-level estimation
remains expensive even with a linear response because solving a large,
ill-conditioned inverse problem is fundamentally different from merely
executing a forward neural-network pass. The low-dimensional gravity model
addresses this boundary through modeling as well as software efficiency.
Section~\ref{sec:why-raptor-scales} explains why the laptop-scale reduced
workflow depends jointly on bounded timetable search, reuse of fixed responses,
and parameter reduction. RAPTOR by itself does not remove the cost of exhaustive
preprocessing or an unnecessarily large response representation.
